% Generated draft: requires mathematical and source review.
\documentclass[11pt]{article}
\usepackage[margin=25mm]{geometry}
\usepackage{fontspec}
\setmainfont{Latin Modern Roman}
\usepackage{amsmath,amssymb,mathtools,unicode-math}
\setmathfont{Latin Modern Math}
\usepackage{enumitem,needspace,xurl,hyperref}
\setlength{\parindent}{0pt}
\setlength{\parskip}{5pt}
\setlength{\emergencystretch}{3em}
\newcommand{\N}{\mathbb{N}}
\newcommand{\Z}{\mathbb{Z}}
\newcommand{\Q}{\mathbb{Q}}
\newcommand{\R}{\mathbb{R}}
\newcommand{\C}{\mathbb{C}}
\newcommand{\F}{\mathbb{F}}
\newcommand{\A}{\mathbb{A}}
\newcommand{\PP}{\mathbb{P}}
\newcommand{\E}{\mathbb{E}}
\newcommand{\eps}{\varepsilon}
\newcommand{\dd}{\,\mathrm{d}}
\begin{document}
\section*{UnsolvedMath ID: 20001657}
\textbf{Source problem code:} catalogue entry
\subsection{Definitions and mathematical statement}
This aggregate record contains Sections 1 and 2 of the workshop list, including remarks and extraction debris, and ends at the heading of Section 3. It must not be reduced to only the first conjecture.

Let $G$ be locally compact with left Haar measure $\mu$ and a length $L$, meaning $L(1)=0$, $L(g^{-1})=L(g)$ and $L(gh)\leq L(g)+L(h)$, with the usual measurability and bounded-on-compact-sets qualifications. For compactly generated groups use a compatible word-length convention. A standard rapid-decay formulation is: some $s>0,C<\infty$ satisfy
\[
 \|\lambda_G(f)\|\leq C\left(\int_G|f(g)|^2(1+L(g))^{2s}\,d\mu(g)\right)^{1/2}
\]
for compactly supported test functions, where $\lambda_G$ is the left regular representation. The source's representation definition has damaged absolute values, powers and quantifiers. Under the standard weighted-$L^2$ interpretation, a nonzero strongly continuous unitary representation $\pi$ has RD when
\[
 \int_G\frac{|\langle\pi(g)v,w\rangle|^2}{(1+L(g))^{2s}}\,d\mu(g)
 \leq C^2\|v\|^2\|w\|^2
\]
for all vectors $v,w$. The exponent/constant normalization above is explicitly fixed, not claimed recovered verbatim from the corrupted display.

\textbf{Section 1: groups and lattices.}
\begin{enumerate}
\item Generalized Valette conjecture: if a locally compact compactly generated group $G$ has RD and $\Gamma\leq G$ is a cocompact lattice, must $\Gamma$ have RD? A lattice is discrete of finite covolume; cocompact means the quotient is compact. The original Valette case takes semisimple Lie $G$. The workshop records positive cases [Ch1] for cocompact lattices in finite products of rank-one simple Lie groups and [La1] for the groups denoted $\operatorname{SL}(3,K)$ with $K=\R,\C,\mathbb H,\mathbb O$, and says higher-rank cases remain open. The quaternionic/octonionic notation requires its source convention. Its example ``$\operatorname{SO}(2,3)=\operatorname{Sp}(4,\R)$'' is not interpreted as a literal equality of global Lie groups; finite-cover/local-isomorphism conventions need checking. It calls the converse an easy result [Jo1]. These are historical source statements, not independently updated status claims.
\item Weaker lattice conjecture: if $\Gamma_1,\Gamma_2$ are cocompact lattices in the same $G$, does RD of $\Gamma_1$ imply RD of $\Gamma_2$? The source also asks for invariance under measurable topological coupling: a locally compact $\Gamma_1\times\Gamma_2$-space with an invariant $\sigma$-compact Borel measure and proper cocompact actions of both groups. The exact meaning of the measure's stated $\sigma$-compact qualification is not clarified.
\item The source notes noncocompact lattices often fail RD because they contain solvable subgroups of exponential growth. This is contextual evidence, not a claim about every nonuniform lattice.
\item Find a useful variation of RD valid for every lattice in connected semisimple Lie groups. Suggested targets are relative RD for $\operatorname{SL}(n,\Z)$, $n\geq3$, relative to a maximal solvable subgroup, and relative RD for solvable groups with respect to exponentially distorted cyclic subgroups. No definition of relative RD or required consequences is supplied. The source says ordinary RD fails for these higher-rank integer lattices because of exponentially growing solvable subgroups.
\item If $H$ is a cocompact normal subgroup of a locally compact compactly generated RD group $G$, must $H$ have RD? Closedness and compatible induced length should be specified. The source says the converse follows by inducing the regular representation of $H$ to $G$.
\item Given $1\to H\to\Gamma\to P\to1$ of finitely generated groups with $P$ having RD, is $\Gamma$ RD exactly when $H$ has RD for the length induced from $\Gamma$? Intrinsic word length on a distorted $H$ cannot be silently substituted. The supplied research summary says Garncarek answered this short-exact-sequence question affirmatively in 2015, distinguishing it from the remaining lattice question.
\item Is every cocompact lattice in a semisimple Lie group $(**)$-relatively hyperbolic with respect to quasi-flats in the sense of [DS2]? The $(**)$ condition and permitted quasi-flat family are not defined in this record; ordinary relative hyperbolicity is not substituted.
\end{enumerate}

\textbf{Section 2: unitary representations and workshop remarks.} The source records that $\lambda_G$ has representation RD exactly when $G$ has group RD; RD passes to weakly contained representations; and a decomposition into irreducibles with uniform RD constants yields RD for the representation with the same constants. A suitable decomposition framework and uniformity are essential, and no existence of such an irreducible decomposition for every arbitrary representation is added.

M. G. Cowling's workshop remark says that if a separated quotient $Q$ of $G$ has a quasi-$G$-invariant measure $\nu$ and the associated unitary representation on $L^2(Q,\nu)$ has RD, then $G$ has RD. In particular, for an amenable closed subgroup $P$, it says $G$ has RD if and only if the representation on $L^2(G/P)$ has RD. The required quotient, measure and unitary Radon--Nikodym conventions are not fully stated. The source further records that the trivial representation has RD exactly for polynomial-growth groups, under its word-length setting. It asks (4) whether a finitely generated group without group RD can nevertheless have a nonzero unitary representation with RD, perhaps by a group extension, and (5) which groups have RD for at least one such representation. Both questions are retained; a zero-dimensional representation is not used to trivialize them.

The supplied aid leaves generalized Valette lattice descent apparently open and gives a conditional coefficient-sampling reduction. For a uniform lattice in a locally compact second-countable RD group, it assumes a uniform polynomial point-to-local-average bound, named (SMV), for coefficients of one fundamental-domain copy of the lattice regular representation. It reports ambient exponent $s$ then gives lattice exponent $s+a$, where $a$ is the sampling-loss exponent, with finite overlap controlled by
\[
 |\Gamma\cap UU^{-1}|
\]
for the fixed relatively compact sampling neighborhood $U$. The exact (SMV) inequality, fundamental-domain embedding, constant and length comparison are not supplied; therefore this is a conditional reduction description, not an unconditional lattice-descent proof. The final Section 3 heading, ``RD for length functions'', contains no substantive question in this record.
\subsection{Short English statement}
This merged workshop extract asks several rapid-decay questions for lattices, extensions, relative formulations and unitary representations. All are retained, including Cowling's remarks and the historical references. The supplied progress solves the extension question in the literature and reduces lattice descent to an extra local coefficient estimate; the missing estimate and damaged definition prevent treating that reduction as a full solution.
\subsection{Sources}
\begin{itemize}
\item Local UnsolvedMath registry, permanent ID 20001657, exact source code catalogue entry; canonical \texttt{aim-geometric-group-theory-notes.json}, \texttt{notes[165]}, source item 1, workshop \emph{The property of rapid decay}; \url{https://aimath.org/WWN/rapiddecay/rapiddecay.pdf}. Supplied verification identifies the eight-page list dated 21 March 2006 and accidental concatenation of Sections 1 and 2, page numbers, and the Section 3 heading, compared with source pages 1--3. That comparison is attributed, not performed here.
\item Original reference labels: [Ch1] for finite products of rank-one simple Lie groups, [La1] for the stated $\operatorname{SL}(3,K)$ cases, [Jo1] for the converse lattice implication, and [DS2] for $(**)$ relative hyperbolicity. No bibliographic expansions are supplied; none is invented. Named attributions include Valette and M. G. Cowling's workshop remark.
\item The supplied summary credits Garncarek (2015) with the affirmative short-exact-sequence result but provides no full citation. The original literature field is empty; the substantial original remarks are retained in the mathematical text.
\item Supplied metadata: PARTIAL-PROGRESS, registry status ``partially\_solved'', statement status ``exact''; attempt 1, result ``partial\_result'', problem status at run ``partially\_solved'', recorded rigor ``proved'', novelty confidence ``low'', candidate type ``reduction''. The machine aid is not peer reviewed and explicitly requires independent verification before citation. No external source was visited.
\end{itemize}
\end{document}
